\section{Theory}
\label{sec:theory}
This section summarizes the signal theory needed for amplitude modulation (AM) and its demodulation methods.

\subsection{Amplitude Modulation Fundamentals}
For standard AM with a large carrier (DSB-LC), the transmitted signal is
\begin{equation}
s(t)=\left[A_c + A_m\,m(t)\right]\cos(2\pi f_c t),
\label{eq:am_general}
\end{equation}
where $m(t)$ is the baseband message (typically normalized so that $|m(t)|\le 1$).

For a single-tone message,
\begin{equation}
m(t)=\cos(2\pi f_m t),
\label{eq:single_tone}
\end{equation}
so
\begin{equation}
s(t)=\left[A_c + A_m\cos(2\pi f_m t)\right]\cos(2\pi f_c t).
\label{eq:am_single_tone}
\end{equation}

The parameters are:
\begin{itemize}
\item $A_c$: carrier amplitude,
\item $A_m$: message scaling amplitude,
\item $f_c$: carrier frequency,
\item $f_m$: message frequency (single-tone case).
\end{itemize}

The modulation index is
\begin{equation}
\mu=\frac{A_m}{A_c}.
\label{eq:mu_def}
\end{equation}
AM works by placing the message in the \emph{envelope} of a high-frequency carrier:
\begin{equation}
e(t)=A_c + A_m\,m(t)=A_c\left[1+\mu\,m(t)\right].
\end{equation}
In DSB-LC, the carrier term is deliberately transmitted, which enables simple envelope detection.

\subsection{Spectral Characteristics of AM}
Let $M(f)$ be the Fourier transform of $m(t)$. From \eqref{eq:am_general},
\begin{equation}
S(f)=\frac{A_c}{2}\!\left[\delta(f-f_c)+\delta(f+f_c)\right]
+\frac{A_m}{2}\!\left[M(f-f_c)+M(f+f_c)\right].
\label{eq:am_spectrum_general}
\end{equation}
This shows two effects:
\begin{itemize}
\item a carrier line at $\pm f_c$,
\item shifted copies of baseband spectrum around $\pm f_c$ (frequency translation).
\end{itemize}

For $m(t)=\cos(2\pi f_m t)$, the sidebands are at
\begin{equation}
f_c-f_m \quad \text{and} \quad f_c+f_m.
\end{equation}
Hence the AM bandwidth is
\begin{equation}
B=2f_m
\end{equation}
for single-tone modulation, and more generally
\begin{equation}
B=2B_m,
\end{equation}
where $B_m$ is the message bandwidth.

\subsection{Coherent Demodulation Theory}
Coherent (synchronous) demodulation multiplies the received AM signal by a local carrier:
\begin{equation}
y(t)=s(t)\cos(2\pi f_c t+\phi),
\label{eq:coherent_multiply}
\end{equation}
where $\phi$ is the receiver phase error.

Using $\cos\alpha\cos\beta=\tfrac{1}{2}[\cos(\alpha-\beta)+\cos(\alpha+\beta)]$,
\begin{equation}
y(t)=\frac{1}{2}\left[A_c+A_m m(t)\right]\cos\phi
+\frac{1}{2}\left[A_c+A_m m(t)\right]\cos(4\pi f_c t+\phi).
\label{eq:coherent_expand}
\end{equation}
Therefore:
\begin{itemize}
\item the first term is the desired baseband term,
\item the second term is centered around $2f_c$ and must be removed.
\end{itemize}

After low-pass filtering:
\begin{equation}
y_{\mathrm{LPF}}(t)=\frac{1}{2}\left[A_c+A_m m(t)\right]\cos\phi.
\label{eq:coherent_lpf}
\end{equation}

Three key consequences follow directly from \eqref{eq:coherent_lpf}:
\begin{itemize}
\item \textbf{Why low-pass filter?} To reject the $2f_c$ mixing product.
\item \textbf{Why remove DC?} The term $\tfrac{1}{2}A_c\cos\phi$ is a DC offset.
\item \textbf{Why scale amplitude?} The message is attenuated by $\tfrac{1}{2}\cos\phi$ (and any LO gain factors), so a gain correction is needed.
\end{itemize}

Coherence is critical: if $\phi\neq 0$, the output is reduced by $\cos\phi$; if frequency is mismatched, an additional time-varying term appears and the recovered baseband is distorted.

\subsection{Envelope Detection Theory}
Envelope detection is a nonlinear alternative to coherent detection. Conceptually it consists of:
\begin{enumerate}
\item rectification,
\item low-pass filtering,
\item DC removal.
\end{enumerate}

For correct operation, the envelope must not invert. With $|m(t)|\le 1$:
\begin{equation}
e_{\min}=A_c-A_m=A_c(1-\mu),
\end{equation}
so the no-inversion condition is
\begin{equation}
\mu \le 1.
\label{eq:no_inversion}
\end{equation}
If $\mu>1$, envelope inversion (over-modulation) occurs, and the rectified/filtered signal no longer tracks the original message linearly.

Compared with coherent detection:
\begin{itemize}
\item coherent detection is linear (with a synchronized carrier) and more accurate,
\item envelope detection is simpler but approximate and more sensitive to amplitude impairments.
\end{itemize}

\subsection{Rectification Loss and Amplitude Scaling}
Rectification is nonlinear, so there is no single impulse response that applies to all inputs. A useful approximation (also the intuition used in the handout Appendix B) is to model half-wave rectification as multiplication by a periodic gating waveform $q(t)$ that keeps positive carrier half-cycles and suppresses negative ones.

For a 50\% duty-cycle gate synchronized to the carrier:
\begin{equation}
q(t)=\frac{1}{2}+\frac{2}{\pi}\!\left[\cos(\omega_c t)+\frac{1}{3}\cos(3\omega_c t)+\frac{1}{5}\cos(5\omega_c t)+\cdots\right].
\label{eq:gate_series}
\end{equation}
Multiplying $s(t)$ by $q(t)$ creates many frequency components. After low-pass filtering, the recovered envelope term is attenuated by an approximately constant factor, commonly expressed as
\begin{equation}
\hat{e}(t)\approx \frac{1}{\pi}\left[A_c + A_m m(t)\right].
\label{eq:pi_loss}
\end{equation}
This explains the practical $\pi$ scaling compensation: multiplying by $\pi$ (before or after rectification, depending on implementation) restores the expected amplitude. The exact factor depends on the exact rectifier model, but the $\pi$ correction captures the dominant loss mechanism.

\subsection{Role of Low-Pass Filtering}
Low-pass filtering is essential in both demodulators:
\begin{itemize}
\item \textbf{Coherent detection:} remove terms around $2f_c$ and keep baseband.
\item \textbf{Envelope detection:} smooth rectifier ripple/harmonics and retain envelope variation.
\end{itemize}

Cutoff selection rule:
\begin{equation}
f_{\mathrm{LP}} \gtrsim B_m \quad \text{and} \quad f_{\mathrm{LP}} \ll f_c.
\end{equation}
So the message passes with low distortion while carrier-related components are rejected.

Filter-order tradeoff:
\begin{itemize}
\item higher order gives stronger stopband rejection,
\item but increases phase delay/transient settling.
\end{itemize}
Butterworth filters are widely used because their passband is maximally flat (no ripple), which helps preserve message amplitude in-band.

\subsection{Noise Effects in AM Systems}
Noise affects AM demodulators differently:
\begin{itemize}
\item \textbf{Envelope detection:} directly tracks amplitude, so amplitude noise appears strongly in the output.
\item \textbf{Coherent detection:} generally more robust because synchronized mixing and LPF provide better selectivity against out-of-band components.
\end{itemize}

The recovered message amplitude grows with $A_m$ (hence with $\mu$), so at very small $\mu$ the desired baseband can be visually dominated by noise. Increasing $\mu$ improves message visibility up to the over-modulation limit $\mu\le 1$ for envelope detection.

Increasing receiver gain raises both desired signal and noise/interference. In practice this can make weaker disturbances visible in time and frequency plots, even when the intended AM signal is still present.
